\begin{lemma}[Step 8: the two cutoff probabilities at a mark]
Let $F$ be a finite-dimensional real inner-product space, write
$\mathcal H=\ell^2(F\times\R)$, let $(\Omega,\mathcal M,\mu)$ be a probability space
and let $(\mathcal G_n)_{n\in\N}$ be a filtration of $\mathcal M$. Let
$z:\N\times\Omega\to\mathcal H$, $S:\N\times\Omega\to\Z$,
$\sigma:\N\times\Omega\to\N$,
$p:\N\times\Omega\to F$, $j\in\N$, $\varsigma:\Omega\to\N$ and
$T,Y:\N\times\Omega\to\R$. Let $a,b,c,k\in\N$ and
$\epsilon,r,K,C',B_0,L_K\in\R$. Write $\sigma_j=\sigma(j,\cdot)$,
$p_j=p(j,\cdot)$, and $\mathcal R_j=\{\omega:t_{\sigma_j(\omega)}=r\}$ for the
event that the $j$-th mark sits at its own level, where $z_n=(x_n,t_n)$.
Assume:
\begin{enumerate}
\item\label{it:mcp-eps} $\epsilon>0$ and $r>0$; $c=2^{k-j}$ and $b=8\cdot2^{k-j}$;
$\epsilon c=r$ and $\epsilon b=8r$; $a\le c$; $0\le a$ as a real; and
$b-a>0$ as a real;
\item\label{it:mcp-K} $2\le K$ and the choice-of-$K$ inequality
$8C'\bigl(L_K+\sqrt{B_0}\bigr)\le K$;
\item\label{it:mcp-lat} $t_n=\epsilon\,S_n$ for all $n,\omega$, and each
$\omega\mapsto t_n(\omega)$ is strongly measurable;
\item\label{it:mcp-st} $\sigma_j$ is a $(\mathcal G_n)$-stopping time, both
$\sigma_j$ and $\varsigma$ are $\mathcal M$-measurable, and
$\sigma_j\le\varsigma$ pointwise;
\item\label{it:mcp-tools} the measurability tools of the ambient argument: for
every $(\mathcal G_n)$-stopping time $\zeta$ and every $n$,
$\{\zeta\le n\}\in\mathcal G_n$; evaluating a measurable family of real
functions at one or at two measurable random indices is measurable; and
$\omega\mapsto g(S_k(\omega))$ is measurable for every $g:\Z\to\R$ and every
$k$;
\item\label{it:mcp-TY} $T_m$ is $\mathcal G_m$-measurable with
$T_m=t_{\sigma_j}$ on $\{\sigma_j\le m\}$, and $Y_m$ is
$\mathcal G_m$-measurable with
\[
  Y_m=r^{-1}\bigl|x_m-x_{\sigma_j}-(t_m-r)p_j\bigr|
  \qquad\text{on }\{\sigma_j\le m\};
\]
\item\label{it:mcp-markov} Markov's inequality in the form: for every
$A\in\mathcal M$, every $\mu$-integrable $f\ge0$ and every $\kappa>0$ with
$f\ge\kappa$ on $A$, one has $\mu(A)\,\kappa\le\int f\,d\mu$;
\item\label{it:mcp-trav} the travel bound of Step 7: the running supremum
\[
  \Upsilon(\omega)=\sup_{\sigma_j(\omega)\le n\le\varsigma(\omega)}
    r^{-1}\bigl|x_n-x_{\sigma_j}-(t_n-t_{\sigma_j})p_j\bigr|
\]
is $\mu$-integrable and $\int\Upsilon\,d\mu\le C'(L_K+\sqrt{B_0})$;
\item\label{it:mcp-claimT} the conclusion of Claim T$'$ for
$\sigma_j,\varsigma$ on the window $(a,b)$, as stated in
\noderef{StoppedLatticeSupersolutionBound}: for all $h,u:\Z\to\R$ and every
$A\in\mathcal M$ with $u\ge0$, with
$\tfrac12(h(i+1)+h(i-1))-h(i)\le-u(i)$ for $a<i<b$, and with
$A\cap\{\sigma_j\le n\}\in\mathcal G_n$ for every $n$, one has for every $n$
\[
  \int\mathbf1_A\Bigl(h\bigl(S_{(\sigma_j\vee n)\wedge\varsigma}\bigr)
    +\sum_{\sigma_j\le m<n\wedge\varsigma}u(S_m)\Bigr)d\mu
  \le\int\mathbf1_A\,h\bigl(S_{\sigma_j}\bigr)\,d\mu .
\]
\end{enumerate}
Then both cutoff probabilities are at most $\tfrac18$:
\begin{equation}\label{eq:mcp-both}
  \mu\bigl(\mathcal R_j\cap\{K\le Y_{\varsigma}\}\bigr)\le\frac18
  \qquad\text{and}\qquad
  \mu\bigl(\mathcal R_j\cap\{8r\le t_{\varsigma}\}\bigr)\le\frac18 ,
\end{equation}
both measures being read as real numbers.
\end{lemma}
\begin{proof}
\emph{Measurability of the mark event.} The map
$\omega\mapsto t_{\sigma_j(\omega)}(\omega)$ is measurable, by
\ref{it:mcp-tools} applied to the measurable index $\sigma_j$ of
\ref{it:mcp-st} and to the family $\omega\mapsto t_n(\omega)$, measurable by
\ref{it:mcp-lat}; so $\mathcal R_j$, its preimage of the single point $r$, lies
in $\mathcal M$. Moreover $\mathcal R_j$ is adapted at $\sigma_j$: for each $n$,
\[
  \mathcal R_j\cap\{\sigma_j\le n\}
   =\bigl(T_n^{-1}\{r\}\bigr)\cap\{\sigma_j\le n\}\in\mathcal G_n ,
\]
because $T_n$ is $\mathcal G_n$-measurable and agrees with $t_{\sigma_j}$ on
$\{\sigma_j\le n\}$ by \ref{it:mcp-TY}, and because
$\{\sigma_j\le n\}\in\mathcal G_n$ by \ref{it:mcp-tools} and
\ref{it:mcp-st}.

\emph{The first cutoff.} Put
$\mathcal B^1=\mathcal R_j\cap\{K\le Y_{\varsigma}\}$, a measurable set by the
previous paragraph together with the measurability of
$\omega\mapsto Y_{\varsigma(\omega)}(\omega)$, which is \ref{it:mcp-tools}
applied to the measurable index $\varsigma$ and the family $Y_m$ of
\ref{it:mcp-TY}. On $\mathcal B^1$ we have $K\le\Upsilon$: indeed
$\sigma_j\le\varsigma$ by \ref{it:mcp-st}, so $\varsigma$ is one of the indices
in the supremum, and at $m=\varsigma$ the summand equals $Y_\varsigma$ by
\ref{it:mcp-TY} --- using $t_{\sigma_j}=r$ on $\mathcal R_j$, which is exactly
what makes the two displayed expressions agree --- so
$Y_\varsigma\le\Upsilon$, and $K\le Y_\varsigma$ on $\mathcal B^1$. The
function $\Upsilon$ is nonnegative, being a supremum of absolute values, and
$\mu$-integrable by \ref{it:mcp-trav}; and $K>0$ by \ref{it:mcp-K}. Markov's
inequality \ref{it:mcp-markov} therefore gives
$\mu(\mathcal B^1)\,K\le\int\Upsilon\,d\mu\le C'(L_K+\sqrt{B_0})$ by
\ref{it:mcp-trav}. Dividing by $K>0$ and using the choice of $K$ from
\ref{it:mcp-K}, which says $8C'\bigl((1+\log(K+10))+\sqrt{B_0}\bigr)\le K$ and
hence a fortiori $8C'(L_K+\sqrt{B_0})\le K$ at the value $L_K$ the travel bound
is stated with,
\[
  \mu(\mathcal B^1)\le\frac{C'(L_K+\sqrt{B_0})}{K}\le\frac18 .
\]

\emph{The second cutoff.} Let
\[
  \ell(i)=\begin{cases}0,&i<a,\\ \dfrac{i-a}{b-a},&a\le i\le b,\\ 1,&i>b,\end{cases}
\]
well defined since $b-a>0$ by \ref{it:mcp-eps}. Then $\ell\ge0$, $\ell$ is
bounded by $1$, and $\ell$ is affine on $\{a,\dots,b\}$, so its second
difference vanishes for $a<i<b$; hence
$\tfrac12(\ell(i+1)+\ell(i-1))-\ell(i)=0\le-u(i)$ there with $u\equiv0$. Apply
\ref{it:mcp-claimT} with $h=\ell$, $u\equiv0$ and $A=\mathcal R_j$, whose two
hypotheses --- $\mathcal R_j\in\mathcal M$ and
$\mathcal R_j\cap\{\sigma_j\le n\}\in\mathcal G_n$ --- were verified above. For
each $n$ the empty-sum term drops and we get
\begin{equation}\label{eq:mcp-claim}
  \int\mathbf1_{\mathcal R_j}\,
    \ell\bigl(S_{(\sigma_j\vee n)\wedge\varsigma}\bigr)d\mu
  \le\int\mathbf1_{\mathcal R_j}\,\ell\bigl(S_{\sigma_j}\bigr)\,d\mu .
\end{equation}

We bound the right side. On $\mathcal R_j$ we have $t_{\sigma_j}=r$, so
$\epsilon S_{\sigma_j}=r=\epsilon c$ by \ref{it:mcp-lat} and
\ref{it:mcp-eps}; cancelling $\epsilon>0$ gives $S_{\sigma_j}=c$. Hence, on
$\mathcal R_j$, using $a\le c\le b$ --- the first from \ref{it:mcp-eps}, the
second because $c\le 8c=b$ in the reals by $\epsilon b=8r=8\epsilon c$ and
$\epsilon>0$ ---
\[
  \ell\bigl(S_{\sigma_j}\bigr)=\frac{c-a}{b-a}\le\frac{c}{b}=\frac18 ,
\]
the inequality because $a\ge0$ makes $(c-a)/(b-a)\le c/b$ whenever
$0\le a\le c\le b$ and $b>a$, and the last step by $b=8c$. Since
$\mathbf1_{\mathcal R_j}\le1$ and the integrand is nonnegative, the right side
of \eqref{eq:mcp-claim} is at most $\tfrac18$.

For the left side, take $n$ large enough that
$(\sigma_j\vee n)\wedge\varsigma=\varsigma$ --- which holds once $n\ge\varsigma$
at the given $\omega$, since then $\sigma_j\vee n\ge n\ge\varsigma$ --- and
observe that on $\mathcal R_j\cap\{8r\le t_{\varsigma}\}$ we have
$\epsilon S_\varsigma=t_\varsigma\ge8r=\epsilon b$, so $S_\varsigma\ge b$ and
$\ell(S_\varsigma)=1$. As $\ell\ge0$ everywhere, this gives the pointwise bound
$\mathbf1_{\mathcal R_j\cap\{8r\le t_\varsigma\}}
 \le\mathbf1_{\mathcal R_j}\,\ell(S_\varsigma)$. Integrating and using
\eqref{eq:mcp-claim} at such $n$ together with the bound on its right side,
\[
  \mu\bigl(\mathcal R_j\cap\{8r\le t_{\varsigma}\}\bigr)
   \le\int\mathbf1_{\mathcal R_j}\,\ell\bigl(S_\varsigma\bigr)\,d\mu
   \le\frac18 .
\]
This is the second half of \eqref{eq:mcp-both}, and the proof is complete.
\end{proof}
